\documentclass[a4paper,11pt]{article}
        \usepackage[top=15mm,bottom=18mm,left=16mm,right=16mm]{geometry}
        \usepackage{fontspec}
        \usepackage{xeCJK}
        \setmainfont{Times New Roman}
        \setCJKmainfont{Yu Gothic}
        \setmonofont{Consolas}
        \setCJKmonofont{Yu Gothic}
        \usepackage{booktabs}
        \usepackage{longtable}
        \usepackage{tabularx}
        \usepackage{array}
        \usepackage{enumitem}
        \usepackage{hyperref}
        \usepackage{listings}
        \usepackage{xcolor}
        \usepackage{amsmath}
        \usepackage{amssymb}
        \hypersetup{
          colorlinks=true,
          linkcolor=blue,
          urlcolor=blue,
          pdftitle={live系 node 構成ビルダ},
          pdfauthor={Codex}
        }
        \lstset{
          basicstyle=\ttfamily\small,
          breaklines=true,
          columns=fullflexible,
          keepspaces=true,
          frame=single,
          framerule=0.2pt,
          rulecolor=\color{black},
          showstringspaces=false
        }
        \setlength{\parskip}{0.42em}
        \setlength{\parindent}{1em}
        \setlength{\emergencystretch}{3em}
        \renewcommand{\arraystretch}{1.15}
        \title{07. live系 node 構成ビルダ}
        \author{solar\_ws0129-main}
        \date{2026-07-29}
        \begin{document}
        \maketitle
        \tableofcontents
        \clearpage

        \section{対象}
        \begin{tabularx}{\linewidth}{>{\raggedright\arraybackslash}p{0.18\linewidth} X}
        \toprule
        項目 & 内容 \\
        \midrule
        ファイル & \path|mpc_solarcar/live_launch.py| \\
        ソースSHA-256 & \path|bee332a416c117f3f953081a7ba7f0c296ce0a3df527510ff7ee140059baceb7| \\
        種別 & Python \\
        区分 & launch helper \\
        役割 & profile を読み、live と live\_wifi で共通に使う node 群を Python から組み立てる。 \\
        起動文脈 & launch/solar\_race\_live*.launch.py から呼ばれる。 \\
        \bottomrule
        \end{tabularx}

        \section{このファイルを読む理由}
        profile を読み、live と live\_wifi で共通に使う node 群を Python から組み立てる。

        \begin{itemize}[leftmargin=1.6em]
        \item forecast CSV の raw/corrected パスを用意する。
\item mpc\_node、logger、dashboard、preflight などの Node action を返す。
\item use\_wifi の有無で planner が読む forecast CSV を切り替える。
        \end{itemize}

        \section{呼び出し関係}
        \subsection{どこから呼ばれるか}
        \begin{itemize}[leftmargin=1.6em]
        \item \path|launch/solar_race_live.launch.py|
\item \path|launch/solar_race_live_wifi.launch.py|
        \end{itemize}

        \subsection{次に読むと理解が進むファイル}
        \begin{itemize}[leftmargin=1.6em]
        \item \path|mpc_solarcar/solar_profile.py|
\item \path|mpc_solarcar/mpc_node.py|
\item \path|mpc_solarcar/speed_command_bridge_node.py|
        \end{itemize}

        \subsection{同一リポジトリ内の主要依存}
        \begin{itemize}[leftmargin=1.6em]
        \item \path|mpc_solarcar/solar_profile.py|
        \end{itemize}

        \section{ソース冒頭の把握}
        モジュール docstring または先頭意図の要約:

        モジュール docstring は明示されていない。

        \subsection{代表コード断片}
        \begin{lstlisting}[language=Python]
if use_wifi and bool(wind_cfg.get("enabled", True)):
        mpc_forecast_csv = corrected_forecast_csv

    nodes = [
        Node(
            package="mpc_solarcar",
            executable="solar_preflight_node",
            name="solar_preflight_node",
            parameters=[
                {
                    "require_speed": True,
                    "require_distance": True,
                    "require_battery": True,
                    "require_planner": True,
                    **preflight_cfg,
\end{lstlisting}


        \section{主要構造の読み取り}
        主要関数は cfg\_path, live\_forecast\_paths, build\_live\_nodes。 launch から起動する Node action は 9 件。

        \subsection{主要クラス・関数}
            \begin{longtable}{p{0.07\linewidth} p{0.16\linewidth} p{0.25\linewidth} p{0.40\linewidth}}
            \toprule
            行 & 種別 & 名前 & 補足 \\
            \midrule
            11 & function & \path|cfg_path| & args: profile\_path, raw \\
15 & function & \path|_drop_keys| & args: payload \\
20 & function & \path|live_forecast_paths| & args: profile\_path, cfg \\
39 & function & \path|build_live_nodes| & args: profile\_path, cfg \\
            \bottomrule
            \end{longtable}

        \subsection{ファイルを上から読んだときの定義順}
            \begin{longtable}{p{0.07\linewidth} p{0.84\linewidth}}
            \toprule
            行 & 内容 \\
            \midrule
            11 & 関数 cfg\_path を定義する。 \\
15 & 関数 \_drop\_keys を定義する。 \\
20 & 関数 live\_forecast\_paths を定義する。 \\
39 & 関数 build\_live\_nodes を定義する。 \\
            \bottomrule
            \end{longtable}

        \subsection{import 群の詳細}
            \begin{longtable}{p{0.07\linewidth} p{0.33\linewidth} p{0.50\linewidth}}
            \toprule
            行 & import 文 & このファイルで読む理由 \\
            \midrule
            1 & \path|from __future__ import annotations| & 型ヒントの遅延評価を有効にし、自己参照型や forward reference を安全に扱うため。 このファイル内での使用位置は少ないか、間接利用である。 \\
3 & \path|import shutil| & shutil モジュールを利用するため。 このファイル内での主な使用位置は L35。 \\
4 & \path|from pathlib import Path| & ファイルやディレクトリを安全に扱うため。 このファイル内での主な使用位置は L32, L33。 \\
6 & \path|from launch_ros.actions import Node| & ROS 2 ノード起動 action を launch から記述するため。 このファイル内での主な使用位置は L58, L72, L102, L113, L131, L140, L142, L145, ...。 \\
8 & \path|from .solar_profile import get_path, get_section, resolve_relative_path| & profile YAML 読込と検証 から get\_path, get\_section, resolve\_relative\_path を読み込み、このファイルの内部処理を分担させるため。 実体ファイルは mpc\_solarcar/solar\_profile.py。 このファイル内での主な使用位置は L12, L21, L22, L23, L24, L40, L41, L42, ...。 \\
            \bottomrule
            \end{longtable}

        \section{このファイルを読む前に必要な基礎知識}
        次の章は、構文やROS用語を既知と仮定しないための説明である。

        \subsection{プログラム、プロセス、メモリ、オブジェクトを区別する}
ソースファイルはディスク上の文字列であり、それ自体は走っていない。Python実行可能プログラムがソースを読み、OSがその実行を一つのプロセスとして管理する。プロセスには仮想メモリ、開いているファイル、スレッド、終了コードなどが対応する。


\begin{lstlisting}
ソースファイル -> Pythonインタプリタが読み込む -> OS上のプロセス -> Pythonオブジェクトをメモリに生成 -> 関数やcallbackを実行
\end{lstlisting}


「メモリ上に生成する」とは、実行中プロセスが使う記憶領域に、その値の型、属性、参照関係を表す実体を用意することである。変数は実体そのものというより、そのオブジェクトを指す名前として理解するとPythonの代入が読みやすい。


\begin{lstlisting}[language=python]
node = MPCNode()
alias = node
# nodeとaliasは同じオブジェクトを参照する。
\end{lstlisting}


一つのプロセスには複数スレッドを持たせられる。スレッドは同じプロセスのメモリを共有するため、受信callbackと最適化callbackが同じself.zを書き換える場合は実行順と排他を検討する必要がある。

\paragraph{根拠資料}
\begin{itemize}[leftmargin=1.6em]
\item \href{https://docs.python.org/3/tutorial/classes.html}{Python公式チュートリアル: Classes}
\end{itemize}

\subsection{名前、代入、型、参照}
Pythonの代入文は、右辺を先に評価し、その結果のオブジェクトへ左辺の名前を結び付ける。`x = f()`では、まずfを呼び、その戻り値が得られてからxが更新される。


`float(...)`、`int(...)`、`bool(...)`は型名を呼び出して値を変換する。外部CSV、YAML、ROS parameterから来た値は型が期待どおりとは限らないため、このリポジトリでは明示変換が多い。


`None`は値が無いことを示す単一のオブジェクト、`math.nan`は浮動小数点値ではあるが有効な数値ではないことを示す。両者は用途が異なるため、`is None`と`math.isfinite`を使い分ける。

\paragraph{根拠資料}
\begin{itemize}[leftmargin=1.6em]
\item \href{https://docs.python.org/3/tutorial/classes.html}{Python公式チュートリアル: Classes}
\end{itemize}

\subsection{関数、引数、戻り値、スコープ}
`def`は関数本体をその場で実行する文ではない。関数オブジェクトを作り、その名前を現在の名前空間へ登録する。`f`は関数そのもの、`f()`はその関数を呼んだ結果である。


仮引数は関数定義側の受け取り名、実引数は呼出側から渡す値である。位置引数、キーワード引数、既定値付き引数、`*args`、`**kwargs`は渡し方が異なる。


\begin{lstlisting}[language=python]
def clip_speed(value, minimum=0.0, maximum=110.0):
    return min(maximum, max(minimum, value))

v = clip_speed(120.0, maximum=90.0)
\end{lstlisting}


関数内で作った通常の名前はローカルスコープに属する。メソッドが`self.z`を更新すればオブジェクトの状態が残るが、単に`z`を更新しただけならその呼出しのローカル値である。

\paragraph{根拠資料}
\begin{itemize}[leftmargin=1.6em]
\item \href{https://docs.python.org/3/tutorial/classes.html}{Python公式チュートリアル: Classes}
\end{itemize}

\subsection{module、package、import、console\_scripts}
Pythonファイルはmoduleとして読み込める。複数moduleをディレクトリにまとめたものがpackageである。`from .model import SolarCarModel`の先頭の点は、現在と同じpackage内のmodel moduleを指す相対importである。


import時にはファイルのトップレベル文が上から一度実行される。`def`や`class`は関数・クラスを登録するが、その本体の通常処理は呼び出すまで走らない。


setuptoolsの`console\_scripts`は、端末で使う実行可能名と`package.module:function`を対応付けるインストール時メタデータである。生成される実行用ラッパーはmoduleをimportし、指定関数を呼び、戻り値を終了コードとして扱う小さな入口である。


\begin{lstlisting}
ROS launchのexecutable名 -> install済みconsole script -> mpc_solarcar.mpc_nodeをimport -> main()を呼ぶ
\end{lstlisting}

\paragraph{根拠資料}
\begin{itemize}[leftmargin=1.6em]
\item \href{https://setuptools.pypa.io/en/latest/userguide/entry_point.html}{setuptools公式: Entry Points / Console Scripts}
\item \href{https://docs.python.org/3/tutorial/classes.html}{Python公式チュートリアル: Classes}
\end{itemize}

\subsection{例外、try/finally、with、資源解放}
例外は通常の戻り値とは別経路で異常を呼出元へ伝える。`try/except`は想定した異常を処理し、`finally`は成功・失敗にかかわらず後始末を行う。


`with open(...) as f:`はcontext managerを使い、ブロックを出るとファイルを閉じる。CSVやログの破損を避けるため、開いた資源の所有者と閉じる場所を明確にする。


`except Exception: pass`はノードを止めない利点がある一方、入力異常を隠して原因追跡を難しくする。安全に関係する値では、少なくとも頻度制限付き警告、異常カウンタ、fallback状態のpublishを検討する。



\subsection{launch Action、Node Action、実行可能名、remapping}
`launch\_ros.actions.Node(...)`はrclpyのNode基底クラスではなく、指定したpackageのexecutableをプロセスとして起動するlaunch Actionである。


`DeclareLaunchArgument`はlaunch実行時に受け取る入力欄を宣言し、`LaunchConfiguration`はその値を後で解決するsubstitutionを表す。`perform(context)`は実行時contextから確定文字列を取り出す。


launchの`name`はNode名override、`parameters`は起動時parameter、`remappings`はNodeやtopicの既定名を別名へ対応付ける。launchはPythonファイルからmainという名前を推測せず、executableとしてインストールされたconsole scriptを起動する。

\paragraph{根拠資料}
\begin{itemize}[leftmargin=1.6em]
\item \href{https://docs.ros.org/en/humble/Tutorials/Beginner-CLI-Tools/Understanding-ROS2-Nodes/Understanding-ROS2-Nodes.html}{ROS 2 Humble公式: Understanding nodes}
\item \href{https://setuptools.pypa.io/en/latest/userguide/entry_point.html}{setuptools公式: Entry Points / Console Scripts}
\end{itemize}

\subsection{warm startは何を保存し、どう効くか}
warm startは前回またはoffline探索で得た入力系列を、次の最適化の初期候補として渡すことである。答えを固定するのではなく、探索開始点と候補libraryの一員を与える。


\[
u^{0,\mathrm{new}}_i=\operatorname{interp}\left(s_i^{\mathrm{new}};\,s_j^{\mathrm{old}},u_j^{\ast,\mathrm{old}}\right)
\]

距離基準計画では、現在より後ろの旧制御点を捨て、新しい絶対距離制御点へ補間してshiftする。初期状態や天候が変わればcost評価は新条件で行われるので、warm startが不適切でも他seed、CEM、安全fallbackが補う設計が必要である。


warm startの効き方は、局所法なら収束先と反復数、CEMなら初期meanまたは候補pool、receding horizonなら前回解の時間・距離shiftとして現れる。どの位置に渡しているかを呼出引数まで追う。

\paragraph{根拠資料}
\begin{itemize}[leftmargin=1.6em]
\item \href{https://docs.scipy.org/doc/scipy/reference/generated/scipy.optimize.minimize.html}{SciPy公式: scipy.optimize.minimize}
\item \href{https://link.springer.com/book/10.1007/978-1-4757-4321-0}{Rubinstein and Kroese: The Cross-Entropy Method}
\end{itemize}

\subsection{天候、route、補間、時刻、単位}
予報は時刻だけの系列か、時刻とroute距離の2次元gridになり得る。現在時刻tと距離sがgrid点の間なら、周囲の値から補間して日射、温度、風を求める。


UTC、現地timezone、naive datetimeを混ぜると数時間のずれが発生する。入力でtimezoneを確定し、内部比較はUTC、表示は現地時刻という役割分担が安全である。


route距離のkm、速度のkm/hとm/s、時間の秒、energyのWhとJを跨ぐため、変換係数3.6、3600、1000の意味を各式で確認する。



\subsection{freshness、filter、guard、fallback、fail-safe}
分散システムでは最後に受け取った値が現在も有効とは限らない。受信時刻とtimeoutからfreshnessを判定し、stale値を計画状態へ無条件に同期しない。


filterはnoiseと一時的な飛び値を抑えるが、遅れを生む。slew limitは指令変化率を制限する。安全guardはsolverのcost罰則とは別に、現在出力へ強制制約を適用する最後の防波堤である。


fallbackは失敗時の代替動作を事前に決める設計である。前回計画保持、物理に基づく決定論的入力、停止、低速制限などから、故障modeごとに選ぶ。fallback発生はstatusとlogへ残し、正常解と区別する。



\subsection{制御、状態、入力、モデル予測制御MPC}
制御対象の内部を表す状態をx、操作入力をu、外乱・予報をwとすると、離散モデルは`x[k+1] = f(x[k], u[k], w[k])`と書ける。ソーラーカーではSoC、電池温度、距離、速度などが状態候補、速度目標や駆動トルクが入力候補になる。


\[
\min_{u_0,\ldots,u_{N-1}} \sum_{k=0}^{N-1}\ell(x_k,u_k,w_k)+V_f(x_N)
\]

MPCは現在状態からNステップ先まで予測し、目的関数と制約を満たす入力系列を求める。ただし実際に適用するのは通常先頭入力だけで、次回は新しい実測状態から再び解く。これがreceding horizonである。


予測モデル、目的関数、制約、ホライズン、solver、初期値のどれかが変わると答えも変わる。「MPCを使う」だけでは仕様は決まらず、これらを単位付きで追う必要がある。

\paragraph{根拠資料}
\begin{itemize}[leftmargin=1.6em]
\item \href{https://docs.scipy.org/doc/scipy/reference/generated/scipy.optimize.minimize.html}{SciPy公式: scipy.optimize.minimize}
\end{itemize}

\subsection{ROS graph、Node、topic、service、action、parameter}
ROS graphは、実行中Nodeと、それらが持つpublisher、subscription、service、actionなどの接続関係である。Pythonクラス、プロセス、ROS graph上のNode名は関連するが同一物ではない。


topicは継続データ向けの非同期一方向publish/subscribe、serviceは短いrequest/response、actionは時間のかかる目標へfeedback、cancel、resultを持たせる。parameterはNodeごとの設定値である。


publisherが送った時点とsubscription callbackが実行される時点は同じとは限らない。通信遅延、QoS queue、executor待ちを挟むため、センサ値にはtimestampとfreshness判定が必要になる。

\paragraph{根拠資料}
\begin{itemize}[leftmargin=1.6em]
\item \href{https://docs.ros.org/en/humble/Tutorials/Beginner-CLI-Tools/Understanding-ROS2-Nodes/Understanding-ROS2-Nodes.html}{ROS 2 Humble公式: Understanding nodes}
\item \href{https://docs.ros.org/en/humble/Concepts/Basic/Interfaces-Topics-Services-Actions.html}{ROS 2 Humble公式: Topics, Services, Actions}
\item \href{https://docs.ros.org/en/humble/Concepts/Basic/About-Parameters.html}{ROS 2 Humble公式: Parameters}
\end{itemize}



        \section{関数・クラスを上から順に解説}
        \subsection{L11 関数 cfg\_path}
            \begin{itemize}[leftmargin=1.6em]
              \item 行範囲: L11--L12
              \item 所属: module直下
              \item 定義: \path|cfg_path(profile_path, raw: str) -> str|
              \item docstring: 明示されていない。
              \item このブロックが直接呼ぶ主な関数・メソッド: resolve\_relative\_path, str, strip
              \item 戻り値の要点: resolve\_relative\_path(profile\_path.parent, raw) if str(raw or '').strip() else ''
              \item 主なローカル代入名: 特になし。
              \item 読み取る主なインスタンス属性: 特になし。
              \item 更新する主なインスタンス属性: 特になし。
              \item 明示的に送出する例外: 明示的raiseはない。
              \item 制御構造の規模: 条件分岐 0、ループ 0、try 0
            \end{itemize}
            \paragraph{この定義を読むためのPython構文}
            \begin{itemize}[leftmargin=1.6em]
            \item `->`以後は戻り値の型注釈であり、実行時の値を自動変換する命令ではない。
            \end{itemize}
            \paragraph{上から順の処理}
            \begin{enumerate}[leftmargin=1.8em]
            \item resolve\_relative\_path(profile\_path.parent, raw) if str(raw or '').strip() else '' を返す。
            \end{enumerate}
            \paragraph{代表コード断片}
            \begin{lstlisting}[language=Python]
def cfg_path(profile_path, raw: str) -> str:
    return resolve_relative_path(profile_path.parent, raw) if str(raw or "").strip() else ""
\end{lstlisting}

\subsection{L15 関数 \_drop\_keys}
            \begin{itemize}[leftmargin=1.6em]
              \item 行範囲: L15--L17
              \item 所属: module直下
              \item 定義: \path|_drop_keys(payload: dict, *keys: str) -> dict|
              \item docstring: 明示されていない。
              \item このブロックが直接呼ぶ主な関数・メソッド: items, set
              \item 戻り値の要点: \{key: value for key, value in (payload or \{\}).items() if key not in blocked\}
              \item 主なローカル代入名: blocked, key, value
              \item 読み取る主なインスタンス属性: 特になし。
              \item 更新する主なインスタンス属性: 特になし。
              \item 明示的に送出する例外: 明示的raiseはない。
              \item 制御構造の規模: 条件分岐 0、ループ 0、try 0
            \end{itemize}
            \paragraph{この定義を読むためのPython構文}
            \begin{itemize}[leftmargin=1.6em]
            \item `->`以後は戻り値の型注釈であり、実行時の値を自動変換する命令ではない。
\item 内包表記は、反復・条件・値生成を一つの式で記述してcollectionを作る。
            \end{itemize}
            \paragraph{上から順の処理}
            \begin{enumerate}[leftmargin=1.8em]
            \item blocked に set(keys) の結果を代入する。
\item \{key: value for key, value in (payload or \{\}).items() if key not in blocked\} を返す。
            \end{enumerate}
            \paragraph{代表コード断片}
            \begin{lstlisting}[language=Python]
def _drop_keys(payload: dict, *keys: str) -> dict:
    blocked = set(keys)
    return {key: value for key, value in (payload or {}).items() if key not in blocked}
\end{lstlisting}

\subsection{L20 関数 live\_forecast\_paths}
            \begin{itemize}[leftmargin=1.6em]
              \item 行範囲: L20--L36
              \item 所属: module直下
              \item 定義: \path|live_forecast_paths(profile_path, cfg)|
              \item docstring: 明示されていない。
              \item このブロックが直接呼ぶ主な関数・メソッド: Path, cfg\_path, copy2, exists, get, get\_path, get\_section, mkdir, str
              \item 戻り値の要点: (base\_path, raw\_path, corrected\_path)
              \item 主なローカル代入名: base\_path, corrected\_path, live\_cfg, raw\_path, runtime\_path, target, weather\_cfg, wind\_cfg
              \item 読み取る主なインスタンス属性: 特になし。
              \item 更新する主なインスタンス属性: 特になし。
              \item 明示的に送出する例外: 明示的raiseはない。
              \item 制御構造の規模: 条件分岐 3、ループ 1、try 0
            \end{itemize}
            \paragraph{この定義を読むためのPython構文}
            \begin{itemize}[leftmargin=1.6em]
            \item 特別な構文上の補足は少ない。
            \end{itemize}
            \paragraph{上から順の処理}
            \begin{enumerate}[leftmargin=1.8em]
            \item live\_cfg に get\_section(cfg, 'live') の結果を代入する。
\item weather\_cfg に get\_section(live\_cfg, 'weather') の結果を代入する。
\item wind\_cfg に get\_section(live\_cfg, 'wind\_model') の結果を代入する。
\item base\_path に get\_path(cfg, profile\_path, 'forecast\_csv') の結果を代入する。
\item raw\_path に cfg\_path(profile\_path, str(weather\_cfg.get('raw\_forecast\_csv', ''))) の結果を代入する。
\item 条件 not raw\_path を判定し、真なら内部処理を行う。
\item   raw\_path に cfg\_path(profile\_path, 'outputs/runtime/live\_forecast\_raw.csv') の結果を代入する。
\item corrected\_path に cfg\_path(profile\_path, str(wind\_cfg.get('corrected\_forecast\_csv', ''))) の結果を代入する。
\item 条件 not corrected\_path を判定し、真なら内部処理を行う。
\item   corrected\_path に cfg\_path(profile\_path, 'outputs/runtime/live\_forecast\_corrected.csv') の結果を代入する。
\item (raw\_path, corrected\_path) を順に走査し、各要素を runtime\_path に入れて処理する。
\item   target に Path(runtime\_path) の結果を代入する。
\item   条件 not target.exists() and Path(base\_path).exists() を判定し、真なら内部処理を行う。
\item     target.parent.mkdir(...) を実行する。
\item     shutil.copy2(...) を実行する。
\item (base\_path, raw\_path, corrected\_path) を返す。
            \end{enumerate}
            \paragraph{代表コード断片}
            \begin{lstlisting}[language=Python]
def live_forecast_paths(profile_path, cfg):
    live_cfg = get_section(cfg, "live")
    weather_cfg = get_section(live_cfg, "weather")
    wind_cfg = get_section(live_cfg, "wind_model")
    base_path = get_path(cfg, profile_path, "forecast_csv")
    raw_path = cfg_path(profile_path, str(weather_cfg.get("raw_forecast_csv", "")))
    if not raw_path:
        raw_path = cfg_path(profile_path, "outputs/runtime/live_forecast_raw.csv")
    corrected_path = cfg_path(profile_path, str(wind_cfg.get("corrected_forecast_csv", "")))
    if not corrected_path:
        corrected_path = cfg_path(profile_path, "outputs/runtime/live_forecast_corrected.csv")
    for runtime_path in (raw_path, corrected_path):
        target = Path(runtime_path)
        if not target.exists() and Path(base_path).exists():
            target.parent.mkdir(parents=True, exist_ok=True)
            shutil.copy2(base_path, target)
    return base_path, raw_path, corrected_path
\end{lstlisting}

\subsection{L39 関数 build\_live\_nodes}
            \begin{itemize}[leftmargin=1.6em]
              \item 行範囲: L39--L175
              \item 所属: module直下
              \item 定義: \path|build_live_nodes(profile_path, cfg, *, use_wifi: bool)|
              \item docstring: 明示されていない。
              \item このブロックが直接呼ぶ主な関数・メソッド: Node, \_drop\_keys, append, bool, cfg\_path, float, get, get\_path, get\_section, int, live\_forecast\_paths, str
              \item 戻り値の要点: nodes
              \item 主なローカル代入名: autocal\_cfg, base\_forecast\_csv, command\_cfg, corrected\_forecast\_csv, distance\_cfg, grade\_cfg, live\_cfg, live\_logging\_cfg, logging\_cfg, mpc\_forecast\_csv, nodes, preflight\_cfg, raw\_forecast\_csv, runtime\_cfg, weather\_cfg, wind\_cfg
              \item 読み取る主なインスタンス属性: 特になし。
              \item 更新する主なインスタンス属性: 特になし。
              \item 明示的に送出する例外: 明示的raiseはない。
              \item 制御構造の規模: 条件分岐 6、ループ 0、try 0
            \end{itemize}
            \paragraph{この定義を読むためのPython構文}
            \begin{itemize}[leftmargin=1.6em]
            \item 特別な構文上の補足は少ない。
            \end{itemize}
            \paragraph{上から順の処理}
            \begin{enumerate}[leftmargin=1.8em]
            \item runtime\_cfg に get\_section(cfg, 'runtime') の結果を代入する。
\item logging\_cfg に get\_section(cfg, 'logging') の結果を代入する。
\item live\_cfg に get\_section(cfg, 'live') の結果を代入する。
\item weather\_cfg に get\_section(live\_cfg, 'weather') の結果を代入する。
\item autocal\_cfg に get\_section(live\_cfg, 'autocal') の結果を代入する。
\item command\_cfg に get\_section(live\_cfg, 'command\_bridge') の結果を代入する。
\item distance\_cfg に get\_section(live\_cfg, 'distance') の結果を代入する。
\item grade\_cfg に get\_section(live\_cfg, 'grade') の結果を代入する。
\item wind\_cfg に get\_section(live\_cfg, 'wind\_model') の結果を代入する。
\item live\_logging\_cfg に get\_section(live\_cfg, 'logging') の結果を代入する。
\item preflight\_cfg に get\_section(live\_cfg, 'preflight') の結果を代入する。
\item (base\_forecast\_csv, raw\_forecast\_csv, corrected\_forecast\_csv) に live\_forecast\_paths(profile\_path, cfg) の結果を代入する。
\item mpc\_forecast\_csv に raw\_forecast\_csv の結果を代入する。
\item 条件 use\_wifi and bool(wind\_cfg.get('enabled', True)) を判定し、真なら内部処理を行う。
\item   mpc\_forecast\_csv に corrected\_forecast\_csv の結果を代入する。
\item nodes に [Node(package='mpc\_solarcar', executable='solar\_preflight\_node', name='solar\_preflight\_node', parameters=[\{'require\_speed': True, 'require\_distance': True, 'require\_battery': True, 'require\_planner': True, **preflight\_cfg\}]), Node(package='mpc\_solarcar', executable='mpc\_node', name='mpc\_node', parameters=[\{'forecast\_csv': mpc\_forecast\_csv, 'route\_profile\_csv': get\_path(cfg, profile\_path, 'route\_profile\_csv'), 'speed\_profile\_csv': get\_path(cfg, profile\_path, 'speed\_profile\_csv'), 'stop\_yaml': get\_path(cfg, profile\_path, 'stop\_yaml'), 'drive\_schedule\_yaml': get\_path(cfg, profile\_path, 'drive\_schedule\_yaml'), 'initial\_upper\_policy\_csv': get\_path(cfg, profile\_path, 'initial\_upper\_policy\_csv'), 'drive\_eff\_map': get\_path(cfg, profile\_path, 'drive\_eff\_map'), 'regen\_eff\_map': get\_path(cfg, profile\_path, 'regen\_eff\_map'), 'rint\_map': get\_path(cfg, profile\_path, 'rint\_map'), 'panel\_eff\_map': get\_path(cfg, profile\_path, 'panel\_eff\_map'), 'mppt\_eff\_map': get\_path(cfg, profile\_path, 'mppt\_eff\_map'), 'drive\_map\_eco': get\_path(cfg, profile\_path, 'drive\_map\_eco'), 'drive\_map\_power': get\_path(cfg, profile\_path, 'drive\_map\_power'), 'regen\_map\_eco': get\_path(cfg, profile\_path, 'regen\_map\_eco'), 'regen\_map\_power': get\_path(cfg, profile\_path, 'regen\_map\_power'), 'ocv\_soc\_map': get\_path(cfg, profile\_path, 'ocv\_soc\_map'), 'params\_yaml': str(profile\_path), 'profile\_runtime\_mode': 'live', 'forecast\_time\_mode': str(live\_cfg.get('forecast\_time\_mode', runtime\_cfg.get('forecast\_time\_mode', 'absolute'))), 'forecast\_time\_tz': str(live\_cfg.get('forecast\_time\_tz', runtime\_cfg.get('forecast\_time\_tz', 'Australia/Darwin'))), 'forecast\_start\_time\_utc': str(get\_section(cfg, 'simulation').get('forecast\_start\_time\_utc', get\_section(cfg, 'simulation').get('start\_utc', '')))\}]), Node(package='mpc\_solarcar', executable='dashboard\_node', name='dashboard\_node', parameters=[\{'host': str(runtime\_cfg.get('dashboard\_host', '0.0.0.0')), 'port': int(runtime\_cfg.get('dashboard\_port', 8080))\}]), Node(package='mpc\_solarcar', executable='solar\_logger\_node', name='solar\_logger\_node', parameters=[\{'log\_dir': cfg\_path(profile\_path, str(logging\_cfg.get('log\_dir', 'outputs/logs'))), 'file\_prefix': str(live\_logging\_cfg.get('file\_prefix', 'solar\_live')), 'log\_rate\_hz': float(logging\_cfg.get('log\_rate\_hz', 2.0)), 'output\_speed\_topic': str(command\_cfg.get('output\_speed\_topic', '/vehicle/speed\_cmd\_kmh')), 'output\_drive\_mode\_topic': str(command\_cfg.get('output\_drive\_mode\_topic', '/vehicle/drive\_mode\_cmd'))\}])] の結果を代入する。
\item 条件 bool(command\_cfg.get('enabled', True)) を判定し、真なら内部処理を行う。
\item   nodes.append(...) を実行する。
\item 条件 bool(live\_cfg.get('use\_distance\_node', True)) を判定し、真なら内部処理を行う。
\item   nodes.append(...) を実行する。
\item 条件 bool(live\_cfg.get('use\_grade\_node', True)) を判定し、真なら内部処理を行う。
\item   nodes.append(...) を実行する。
\item 条件 bool(weather\_cfg.get('enabled', True)) を判定し、真なら内部処理を行う。
\item   nodes.append(...) を実行する。
\item 条件 bool(autocal\_cfg.get('enabled', True)) を判定し、真なら内部処理を行う。
\item   nodes.append(...) を実行する。
\item nodes を返す。
            \end{enumerate}
            \paragraph{代表コード断片}
            \begin{lstlisting}[language=Python]
def build_live_nodes(profile_path, cfg, *, use_wifi: bool):
    runtime_cfg = get_section(cfg, "runtime")
    logging_cfg = get_section(cfg, "logging")
    live_cfg = get_section(cfg, "live")
    weather_cfg = get_section(live_cfg, "weather")
    autocal_cfg = get_section(live_cfg, "autocal")
    command_cfg = get_section(live_cfg, "command_bridge")
    distance_cfg = get_section(live_cfg, "distance")
    grade_cfg = get_section(live_cfg, "grade")
    wind_cfg = get_section(live_cfg, "wind_model")
    live_logging_cfg = get_section(live_cfg, "logging")
    preflight_cfg = get_section(live_cfg, "preflight")

    base_forecast_csv, raw_forecast_csv, corrected_forecast_csv = live_forecast_paths(profile_path, cfg)
    mpc_forecast_csv = raw_forecast_csv
    if use_wifi and bool(wind_cfg.get("enabled", True)):
        mpc_forecast_csv = corrected_forecast_csv

    nodes = [
        Node(
            package="mpc_solarcar",
            executable="solar_preflight_node",
            name="solar_preflight_node",
            parameters=[
                {
                    "require_speed": True,
                    "require_distance": True,
                    "require_battery": True,
                    "require_planner": True,
                    **preflight_cfg,
                }
            ],
        ),
        Node(
            package="mpc_solarcar",
...
\end{lstlisting}





        \subsection{launch から起動するノード}
            \begin{longtable}{p{0.07\linewidth} p{0.30\linewidth} p{0.50\linewidth}}
            \toprule
            行 & executable & package / name \\
            \midrule
            58 & \path|solar_preflight_node| & package=mpc\_solarcar, name=solar\_preflight\_node \\
72 & \path|mpc_node| & package=mpc\_solarcar, name=mpc\_node \\
102 & \path|dashboard_node| & package=mpc\_solarcar, name=dashboard\_node \\
113 & \path|solar_logger_node| & package=mpc\_solarcar, name=solar\_logger\_node \\
131 & \path|speed_command_bridge_node| & package=mpc\_solarcar, name=speed\_command\_bridge\_node \\
140 & \path|distance_node| & package=mpc\_solarcar, name=distance\_node \\
142 & \path|grade_node| & package=mpc\_solarcar, name=grade\_node \\
145 & \path|weather_fetch_node| & package=mpc\_solarcar, name=weather\_fetch\_node \\
168 & \path|solar_autocal_node| & package=mpc\_solarcar, name=solar\_autocal\_node \\
            \bottomrule
            \end{longtable}





        \section{処理の流れ}
        \begin{enumerate}[leftmargin=1.7em]
        \item ソース中の主要関数を通じて処理を進める。
        \end{enumerate}

        \section{このファイルの位置づけ}
        この資料群では、\path|mpc_solarcar/live_launch.py| を
        \textbf{launch helper}の中核として扱う。
        したがって、ここで扱う処理は単独で完結せず、
        前段の profile / launch / telemetry と後段の planner / logger / report へ繋がる。

        \end{document}
